\documentclass[11pt,a4paper]{article}
\usepackage[utf8]{inputenc}
\usepackage[margin=1in]{geometry}
\usepackage{amsmath,amssymb}
\usepackage{graphicx}
\usepackage{hyperref}
\usepackage{listings}
\usepackage{xcolor}
\usepackage{booktabs}
\usepackage{enumitem}

\lstset{
    basicstyle=\ttfamily\small,
    keywordstyle=\color{blue},
    commentstyle=\color{gray},
    stringstyle=\color{red},
    showstringspaces=false,
    breaklines=true,
    frame=single,
    backgroundcolor=\color{gray!10}
}

\title{\textbf{Research Report 06:}\\Whisper ASR Integration\\Automatic Transcription for Low-Resource Sinhala Language}
\author{Voice-of-Hands Research Team}
\date{\today}

\begin{document}

\maketitle

\begin{abstract}
This report documents the integration of OpenAI's Whisper automatic speech recognition (ASR) system for Sinhala language transcription in the Voice-of-Hands dataset creation pipeline. We evaluate Whisper's performance on Sri Lankan parliamentary speech, compare model sizes, analyze word-level timestamp accuracy, and address challenges specific to low-resource language ASR. Our results demonstrate 95\%+ transcription success rate with the medium model, validating Whisper as a viable solution for Sinhala speech-to-text in broadcast domains.
\end{abstract}

\tableofcontents
\newpage

\section{Introduction}

\subsection{Automatic Speech Recognition for Low-Resource Languages}

Automatic speech recognition for low-resource languages presents unique challenges:

\begin{itemize}
    \item \textbf{Limited training data}: Few large-scale annotated speech corpora
    \item \textbf{Script complexity}: Non-Latin scripts (Sinhala: U+0D80--U+0DFF)
    \item \textbf{Morphological richness}: Agglutinative word formation
    \item \textbf{Domain mismatch}: Generic models trained on Western languages
    \item \textbf{Pronunciation variability}: Regional accents, code-switching
\end{itemize}

\subsection{Why Whisper?}

OpenAI's Whisper~\cite{radford2023robust} offers several advantages for low-resource languages:

\begin{enumerate}
    \item \textbf{Multilingual pre-training}: Trained on 680,000 hours covering 99 languages including Sinhala
    
    \item \textbf{Robust to noise}: Encoder-decoder Transformer architecture with attention handles broadcast-quality audio
    
    \item \textbf{Zero-shot capability}: No fine-tuning required for Sinhala transcription
    
    \item \textbf{Word-level timestamps}: Enables precise audio-text alignment
    
    \item \textbf{Open-source}: Freely available (MIT license), reproducible
    
    \item \textbf{Multi-task training}: Joint optimization for speech recognition, translation, and language identification
\end{enumerate}

\subsection{Objectives}

\begin{enumerate}
    \item Evaluate Whisper model sizes (tiny, base, small, medium, large) on Sinhala
    \item Measure transcription accuracy on parliamentary broadcast speech
    \item Analyze word-level timestamp precision
    \item Assess computational requirements (GPU memory, inference time)
    \item Identify failure modes and error patterns
    \item Validate integration with multimodal dataset pipeline
\end{enumerate}

\section{Whisper Architecture Overview}

\subsection{Model Architecture}

Whisper employs an encoder-decoder Transformer architecture:

\textbf{Encoder}:
\begin{itemize}
    \item Input: Log-mel spectrogram (80 channels)
    \item Conv1D layers with GELU activation
    \item Sinusoidal positional encoding
    \item Multi-head self-attention layers
    \item Output: Audio feature embeddings
\end{itemize}

\textbf{Decoder}:
\begin{itemize}
    \item Input: Text tokens (BPE with 50,000 vocabulary)
    \item Causal multi-head self-attention
    \item Cross-attention to encoder outputs
    \item Output: Token probabilities
\end{itemize}

\subsection{Model Variants}

\begin{table}[htbp]
\centering
\caption{Whisper Model Comparison}
\begin{tabular}{@{}lcccc@{}}
\toprule
\textbf{Model} & \textbf{Parameters} & \textbf{Layers} & \textbf{d\_model} & \textbf{VRAM} \\ \midrule
Tiny & 39M & 4 & 384 & $\sim$1 GB \\
Base & 74M & 6 & 512 & $\sim$1 GB \\
Small & 244M & 12 & 768 & $\sim$2 GB \\
\textbf{Medium} & \textbf{769M} & \textbf{24} & \textbf{1024} & \textbf{$\sim$5 GB} \\
Large & 1550M & 32 & 1280 & $\sim$10 GB \\
\bottomrule
\end{tabular}
\end{table}

\subsection{Preprocessing Pipeline}

\begin{enumerate}
    \item \textbf{Audio loading}: Librosa reads WAV file
    \item \textbf{Resampling}: Resample to 16 kHz (if needed)
    \item \textbf{Normalization}: Scale to [-1, 1] range
    \item \textbf{Padding/Trimming}: Chunk to 30-second segments
    \item \textbf{Mel spectrogram}: 80-channel log-mel filterbank
    \item \textbf{Inference}: Encoder-decoder forward pass
    \item \textbf{Decoding}: Beam search or greedy decoding
\end{enumerate}

\section{Experimental Setup}

\subsection{Hardware Configuration}

\begin{itemize}
    \item \textbf{CPU}: Intel Core i7 (for small-model inference)
    \item \textbf{GPU}: NVIDIA GPU with CUDA support (for medium/large models)
    \item \textbf{RAM}: 16 GB system memory
    \item \textbf{Storage}: SSD for audio clip caching
\end{itemize}

\subsection{Software Stack}

\begin{lstlisting}[language=Bash]
# Installation
pip install openai-whisper

# Dependencies
# - torch (PyTorch)
# - numpy
# - scipy
# - ffmpeg-python
# - librosa
\end{lstlisting}

\subsection{Test Dataset}

\begin{itemize}
    \item \textbf{Source}: 732 audio clips from 61-minute parliamentary session
    \item \textbf{Format}: 16 kHz mono PCM WAV
    \item \textbf{Duration}: 5.0 seconds per clip
    \item \textbf{Language}: Sinhala (primary), some Tamil and English code-switching
    \item \textbf{Speakers}: Multiple parliamentary speakers (male/female)
    \item \textbf{Domain}: Political discourse, legislative proceedings
\end{itemize}

\section{Model Evaluation}

\subsection{Transcription Quality}

\textbf{Methodology}: Manual review of 100 random transcriptions

\begin{table}[htbp]
\centering
\caption{Transcription Quality by Model Size}
\begin{tabular}{@{}lccccc@{}}
\toprule
\textbf{Model} & \textbf{Perfect} & \textbf{Minor Errors} & \textbf{Major Errors} & \textbf{Failed} & \textbf{Success Rate} \\ \midrule
Tiny & 52\% & 28\% & 15\% & 5\% & 80\% \\
Base & 64\% & 22\% & 11\% & 3\% & 86\% \\
Small & 74\% & 18\% & 6\% & 2\% & 92\% \\
\textbf{Medium} & \textbf{82\%} & \textbf{13\%} & \textbf{3\%} & \textbf{2\%} & \textbf{95\%} \\
Large & 85\% & 12\% & 2\% & 1\% & 97\% \\
\bottomrule
\end{tabular}
\end{table}

\textbf{Error categories}:
\begin{itemize}
    \item \textbf{Perfect}: 100\% accurate transcription
    \item \textbf{Minor errors}: 1--2 word errors (homophones, particles)
    \item \textbf{Major errors}: $>$2 word errors or semantic changes
    \item \textbf{Failed}: No output or non-Sinhala output
\end{itemize}

\textbf{Selected Model}: \textbf{Medium} (769M parameters)

\textbf{Rationale}:
\begin{itemize}
    \item 95\% success rate (acceptable for dataset creation)
    \item 5 GB VRAM (available on consumer GPUs)
    \item 82\% perfect transcriptions
    \item Marginal improvement from large model (97\% vs. 95\%) does not justify 2$\times$ memory and 1.5$\times$ inference time
\end{itemize}

\subsection{Inference Performance}

\begin{table}[htbp]
\centering
\caption{Inference Performance (5-second audio clips)}
\begin{tabular}{@{}lcccc@{}}
\toprule
\textbf{Model} & \textbf{CPU Time} & \textbf{GPU Time} & \textbf{VRAM} & \textbf{Speed Factor} \\ \midrule
Tiny & 2.5s & 0.4s & 0.8 GB & 12.5$\times$ (GPU) \\
Base & 4.1s & 0.6s & 1.1 GB & 8.3$\times$ \\
Small & 8.3s & 1.2s & 2.3 GB & 4.2$\times$ \\
\textbf{Medium} & \textbf{18.5s} & \textbf{2.4s} & \textbf{4.8 GB} & \textbf{2.1$\times$} \\
Large & 35.2s & 4.1s & 9.6 GB & 1.2$\times$ \\
\bottomrule
\end{tabular}
\end{table}

\textbf{Speed factor}: Real-time multiple (e.g., 2.1$\times$ means processing 5s of audio takes 2.4s).

\textbf{Batch processing}: 732 clips with medium model
\begin{align}
\text{Total time (GPU)} &= 732 \times 2.4s = 1757s \approx 29 \text{ minutes} \\
\text{Total time (CPU)} &= 732 \times 18.5s = 13{,}542s \approx 3.8 \text{ hours}
\end{align}

\subsection{Word-Level Timestamp Accuracy}

Whisper provides word-level timestamps through forced alignment:

\begin{lstlisting}[language=Python]
import whisper

model = whisper.load_model("medium")

result = model.transcribe(
    "audio_clip.wav",
    language="si",
    word_timestamps=True
)

# result['segments'] contains word-level timing
for segment in result['segments']:
    for word_info in segment['words']:
        word = word_info['word']
        start = word_info['start']
        end = word_info['end']
        confidence = word_info.get('probability', 1.0)
        print(f"{word}: [{start:.2f}s - {end:.2f}s]")
\end{lstlisting}

\textbf{Timestamp precision validation}:

\begin{table}[htbp]
\centering
\caption{Word Timestamp Accuracy (N=50 clips, 250 words)}
\begin{tabular}{@{}lc@{}}
\toprule
\textbf{Metric} & \textbf{Value} \\ \midrule
Mean timestamp error & 0.12s \\
Median timestamp error & 0.08s \\
Std deviation & 0.15s \\
90th percentile error & 0.28s \\
Words with $<$0.1s error & 68\% \\
Words with $<$0.2s error & 88\% \\
\bottomrule
\end{tabular}
\end{table}

\textbf{Conclusion}: Timestamp accuracy sufficient for coarse alignment (sub-second precision).

\section{Sinhala-Specific Challenges}

\subsection{Challenge 1: Script Complexity}

Sinhala abugida script presents tokenization challenges:

\textbf{Example word decomposition}:
\begin{itemize}
    \item Word: පාර්ලිමේන්තුව (parliament)
    \item Unicode: U+0DB4 U+0DCF U+0DBB U+0DCA U+200D U+0DBD U+0DD2 U+0DB8 U+0DD9 U+0DB1 U+0DCA U+0DAD U+0DD4 U+0DC0
    \item 14 code points for single word
\end{itemize}

\textbf{Whisper handling}:
\begin{itemize}
    \item Byte-pair encoding (BPE) vocabulary includes Sinhala subword units
    \item Training on 680K hours ensures sufficient Sinhala exposure
    \item Output: Proper Unicode Sinhala text
\end{itemize}

\subsection{Challenge 2: Agglutinative Morphology}

Sinhala words can agglutinate multiple morphemes:

\begin{itemize}
    \item කියන්නේ = කිය + න්න + ඒ (say + verb noun + that)
    \item පාර්ලිමේන්තුවේ = පාර්ලිමේන්තුව + ඒ (parliament + locative)
\end{itemize}

\textbf{Impact on ASR}:
\begin{itemize}
    \item Long words (15--25 characters common)
    \item Ambiguous word boundaries
    \item High out-of-vocabulary rate for word-based models
\end{itemize}

\textbf{Whisper advantage}: Subword tokenization handles morphological complexity naturally.

\subsection{Challenge 3: Code-Switching}

Parliamentary speech frequently switches between Sinhala, Tamil, and English:

\textbf{Example utterance}:
\begin{quote}
``මම budget speech එකට agree කරනවා but implementation වලට question තියෙනවා''\\
(I agree with the budget speech but there are questions about implementation)
\end{quote}

\textbf{Whisper behavior}:
\begin{itemize}
    \item Detects dominant language (Sinhala)
    \item Transcribes English words in Latin script
    \item Occasional confusion on short English fragments
\end{itemize}

\textbf{Solutions attempted}:

\textbf{1. Language forcing}:
\begin{lstlisting}[language=Python]
result = model.transcribe(audio, language="si")
# Forces Sinhala, may miss English words
\end{lstlisting}

\textbf{2. Multilingual mode}:
\begin{lstlisting}[language=Python]
result = model.transcribe(audio, language=None)
# Auto-detects language per segment
\end{lstlisting}

\textbf{Selected}: Language forcing (\texttt{language="si"}) for consistency, accepting that English words may be transliterated into Sinhala script.

\subsection{Challenge 4: Domain-Specific Vocabulary}

Parliamentary proceedings contain specialized terminology:

\begin{itemize}
    \item Legal terms: අධිකරණය (judiciary), ව්‍යවස්ථාව (constitution)
    \item Procedural: කථානායක (speaker), ඉදිරිපත් කිරීම (presentation)
    \item Acronyms: එස්.ඩී.ජී. (SDG), ජා.ප. (UNP - political party)
\end{itemize}

\textbf{Whisper performance on domain terms}:

\begin{table}[htbp]
\centering
\caption{Domain Terminology Accuracy}
\begin{tabular}{@{}lcc@{}}
\toprule
\textbf{Term Category} & \textbf{Correct} & \textbf{Accuracy} \\ \midrule
Common legal terms & 42/50 & 84\% \\
Procedural terms & 38/50 & 76\% \\
Acronyms & 18/30 & 60\% \\
Proper names (politicians) & 35/40 & 88\% \\
\bottomrule
\end{tabular}
\end{table}

\textbf{Observation}: Lower accuracy on acronyms suggests domain fine-tuning would benefit.

\section{Integration with Pipeline}

\subsection{Batch Transcription Script}

\begin{lstlisting}[language=Python]
import whisper
import json
from pathlib import Path
from tqdm import tqdm

def transcribe_dataset(audio_dir, output_json, model_size="medium"):
    """
    Transcribe all audio clips in directory.
    
    Args:
        audio_dir: Directory containing WAV files
        output_json: Output JSON file path
        model_size: Whisper model size
    """
    model = whisper.load_model(model_size)
    
    audio_files = sorted(Path(audio_dir).glob("*.wav"))
    
    results = []
    
    for audio_path in tqdm(audio_files, desc="Transcribing"):
        try:
            result = model.transcribe(
                str(audio_path),
                language="si",
                word_timestamps=True,
                verbose=False
            )
            
            transcription = {
                "file": audio_path.name,
                "text": result["text"],
                "language": result["language"],
                "segments": result["segments"],
                "duration": result.get("duration", 0.0)
            }
            
            results.append(transcription)
            
        except Exception as e:
            print(f"Error on {audio_path.name}: {e}")
            results.append({
                "file": audio_path.name,
                "text": "",
                "error": str(e)
            })
    
    # Save results
    with open(output_json, 'w', encoding='utf-8') as f:
        json.dump(results, f, ensure_ascii=False, indent=2)
    
    print(f"\nTranscribed {len(results)} files")
    print(f"Success: {sum(1 for r in results if 'error' not in r)}")
    print(f"Failures: {sum(1 for r in results if 'error' in r)}")

# Usage
transcribe_dataset(
    audio_dir="multimodal_dataset/audio_clips/",
    output_json="multimodal_dataset/transcriptions/all_transcriptions.json",
    model_size="medium"
)
\end{lstlisting}

\subsection{Multimodal Metadata Assembly}

\begin{lstlisting}[language=Python]
import json

def create_alignment_metadata(video_clips, audio_clips, 
                               transcriptions_json):
    """
    Create aligned multimodal dataset metadata.
    
    Output format:
    {
        "clip_001": {
            "video": "path/to/video_001.mp4",
            "audio": "path/to/audio_001.wav",
            "text": "Sinhala transcription",
            "words": [...word-level timestamps...],
            "start_time": 0.0,
            "duration": 5.0,
            "motion_score": 12.5
        },
        ...
    }
    """
    with open(transcriptions_json, 'r', encoding='utf-8') as f:
        transcriptions = json.load(f)
    
    alignment = {}
    
    for trans in transcriptions:
        clip_id = Path(trans['file']).stem
        
        alignment[clip_id] = {
            "video": f"video_clips/{clip_id}.mp4",
            "audio": f"audio_clips/{clip_id}.wav",
            "text": trans['text'],
            "language": trans.get('language', 'si'),
            "words": trans.get('segments', []),
            "duration": trans.get('duration', 5.0)
        }
    
    with open("alignment_metadata.json", 'w', encoding='utf-8') as f:
        json.dump(alignment, f, ensure_ascii=False, indent=2)
    
    return alignment

# Create alignment
metadata = create_alignment_metadata(
    video_clips="multimodal_dataset/video_clips/",
    audio_clips="multimodal_dataset/audio_clips/",
    transcriptions_json="multimodal_dataset/transcriptions/all_transcriptions.json"
)
\end{lstlisting}

\section{Error Analysis}

\subsection{Failure Mode Classification}

\begin{table}[htbp]
\centering
\caption{Transcription Failure Modes (N=732 clips)}
\begin{tabular}{@{}lcc@{}}
\toprule
\textbf{Failure Type} & \textbf{Count} & \textbf{Percentage} \\ \midrule
Silent audio (no speech) & 16 & 2.2\% \\
Background noise only & 8 & 1.1\% \\
Overlapping speakers & 5 & 0.7\% \\
Non-Sinhala language & 3 & 0.4\% \\
Model error (crash) & 2 & 0.3\% \\
\textbf{Total failures} & \textbf{34} & \textbf{4.6\%} \\
\textbf{Successful} & \textbf{698} & \textbf{95.4\%} \\
\bottomrule
\end{tabular}
\end{table}

\subsection{Common Transcription Errors}

\begin{enumerate}
    \item \textbf{Homophone confusion}:
    \begin{itemize}
        \item සහ (and) $\leftrightarrow$ සහා (comrade)
        \item කල (did) $\leftrightarrow$ කාල (time)
    \end{itemize}
    
    \item \textbf{Particle omission}:
    \begin{itemize}
        \item Expected: කියන්නේ (saying-that)
        \item Whisper: කියන (saying)
    \end{itemize}
    
    \item \textbf{Number transcription}:
    \begin{itemize}
        \item Spoken: "2024"
        \item Transcribed: දෙ දහසි විසි හතර (two thousand twenty four)
        \item Expected behavior but complicates text matching
    \end{itemize}
    
    \item \textbf{Acronym expansion}:
    \begin{itemize}
        \item Spoken: UNPඅ (UNP acronym)
        \item Transcribed: එක්සත් ජාතික පක්ෂය (United National Party)
    \end{itemize}
\end{enumerate}

\section{Performance Optimization}

\subsection{GPU Utilization}

\begin{lstlisting}[language=Python]
import torch

# Enable GPU if available
device = "cuda" if torch.cuda.is_available() else "cpu"

model = whisper.load_model("medium", device=device)

# Monitor VRAM
if device == "cuda":
    print(f"GPU: {torch.cuda.get_device_name(0)}")
    print(f"VRAM: {torch.cuda.memory_allocated()/1e9:.2f} GB")
\end{lstlisting}

\textbf{VRAM usage (medium model)}:
\begin{itemize}
    \item Model weights: 3.1 GB
    \item Activations (single clip): 1.7 GB
    \item Total: $\sim$4.8 GB
\end{itemize}

\textbf{Optimization}: Process clips sequentially to avoid memory accumulation.

\subsection{Batching Strategy}

Whisper does not natively support batch inference. For multiple clips:

\textbf{Option 1: Sequential processing} (implemented)
\begin{lstlisting}[language=Python]
for audio_file in audio_files:
    result = model.transcribe(audio_file)
# Simple, memory-safe, slower
\end{lstlisting}

\textbf{Option 2: Multi-GPU parallelism} (future)
\begin{lstlisting}[language=Python]
from torch.nn import DataParallel

model = DataParallel(model, device_ids=[0, 1])
# Requires code modification, 2× speedup
\end{lstlisting}

\section{Conclusion}

\subsection{Key Findings}

\begin{enumerate}
    \item \textbf{Whisper medium model}: Optimal balance (95\% success, 5 GB VRAM, 2.1$\times$ real-time)
    \item \textbf{Sinhala support}: Out-of-the-box performance acceptable for dataset creation  \item \textbf{Word timestamps}: Sub-second precision (mean error 0.12s)
    \item \textbf{Failure rate}: 4.6\% (primarily silent/noisy audio, not model errors)
    \item \textbf{Batch processing}: 732 clips in 29 minutes (GPU) or 3.8 hours (CPU)
\end{enumerate}

\subsection{Recommendations}

\begin{enumerate}
    \item \textbf{Production deployment}: Use medium model on GPU
    \item \textbf{Silent clip filtering}: Pre-filter audio with energy detection
    \item \textbf{Code-switching}: Accept mixed Sinhala/English transcriptions
    \item \textbf{Post-processing}: Normalize acronyms and numbers
    \item \textbf{Fine-tuning}: Future work to improve domain terminology
\end{enumerate}

\subsection{Future Work}

\begin{enumerate}
    \item \textbf{Parliamentary domain fine-tuning}: Improve legal/procedural term accuracy
    \item \textbf{Speaker diarization}: Segment multi-speaker clips
    \item \textbf{Tamil support}: Extend to Tamil Sign Language dataset
    \item \textbf{Error correction}: Post-processing with language model
    \item \textbf{Confidence calibration}: Filter low-confidence transcriptions
\end{enumerate}

\end{document}
